%% ch13 --- Wahadła i planety
%%
%% Napisany we wrześniu 2026. Każda liczba, której czytelnik nie policzy w
%% pamięci, pochodzi z code/measure/ch13.py, który importuje TE SAME funkcje,
%% które drukują listingi (code/ch13/). Wahadło podwójne woła sin i cos, których
%% ostatnia cyfra dwójkowa może się różnić między maszynami, więc nic nie jest
%% zapisywane z późnej części chaotycznego przebiegu: chwila rozejścia jest
%% brana długo przed tym, zanim różnica ostatniej cyfry mogłaby mieć znaczenie,
%% a skrypt odmawia zapisania każdej liczby, którą zmieniłaby inna długość kroku
%% albo puszczenie przesunięte o 1e-12 rad.
%%
%% Bliźniak: chapters/en/ch13-mechanics.tex. Podrozdział za podrozdziałem,
%% makro za makrem, liczba za liczbą; tools/parity.py je porównuje.

\chapter{Wahadła i planety}
\label{ch:mechanics}
\index{wahadło podwójne}\index{Układ Słoneczny}

\noindent
Zegar stojący odmierza czas, bo wahadło to najbardziej regularna rzecz, jaką
większość ludzi kiedykolwiek będzie mieć w domu: puść je, a będzie się kołysać
w tym samym rytmie, dopóki zegar jest nakręcony. Rozdział~\ref{ch:motion}
pokazał dlaczego: stan wahadła to dwie liczby, a bez tarcia krążą one po tej
samej zamkniętej pętli bez końca. Wszechświat jak mechanizm zegara z
rozdziału~\ref{ch:models} zbudowano właśnie na takiej mechanice.

Teraz zawieś drugie wahadło u dołu pierwszego, na zawiasie w miejscu, gdzie
się łączą, i puść je z wysoka. Reguła wciąż pochodzi od Newtona, wciąż jest
dokładna i nie ma w niej nic losowego. Co ono zrobi? Ten rozdział odpowiada
pomiarem, a potem idzie za tym samym pytaniem od wahadła na biurku aż do
planet krążących wokół Słońca, gdzie zadano je po raz pierwszy.

\begin{outcomes}
\outcome{przeprowadzić wahadło podwójne krok po kroku metodą RK4 i odczytać
  jego stan z wydruku}
\outcome{sprawdzić symulację układu bez tarcia, obserwując jego energię, i
  powiedzieć, czego ten sprawdzian dowodzi, a czego nie}
\outcome{zmierzyć, jak daleko od siebie odchodzą dwa niemal identyczne
  puszczenia, i sprawdzić, czy wynik zależy od długości kroku}
\outcome{odróżnić łagodne puszczenie od dzikiego i powiedzieć, co zachowanie
  energii przesądza o przyszłości układu, a czego nie}
\outcome{oddzielić to, co zmierzono na temat chaosu w Układzie Słonecznym, od
  tego, co jest legendą}
\end{outcomes}

\section{Wahadło na wahadle}
\label{sec:ch13-double}
\index{wahadło podwójne!stan}

\noindent
Weź dwa pręty o długości metra i przymocuj drugi zawiasem do końca
pierwszego. Na zawiasie umieść ciężarek o masie jednego kilograma, a drugi
taki sam na wolnym końcu; całą masę niosą ciężarki, a pręty traktujemy jak
nieważkie. Unieś oba ramiona tak, by były odchylone o 120° od pionu, i puść.

\begin{think}
Jedno wahadło kołysze się z doskonałą regularnością. Czy spodziewasz się, że
dwa, jedno zawieszone na drugim, ułożą się w jakiś powtarzający się wzór
\dash{} bardziej skomplikowany niż u jednego wahadła, ale mimo to regularny?
\end{think}
\reveal{Nie. Puszczone z wysoka wahadło podwójne nie układa się w żaden
powtarzający się wzór, jaki dałoby się znaleźć: dolne ramię śmiga dookoła,
czasem przelatując nad samą górą, a ruch wciąż się zmienia. Nic losowego tu
nie działa.}
Różnica tkwi w tym, ile jest liczb. Stan jednego wahadła to dwie liczby: kąt i
to, jak szybko ten kąt się zmienia. Stan wahadła podwójnego to cztery liczby,
dwa kąty i dwie prędkości ich zmian, a rozdział~\ref{ch:motion} mówił, że
gładka reguła potrzebuje co najmniej trzech liczb, żeby mogła być chaotyczna.

Prawa Newtona dają regułę zmiany dla tych czterech liczb: znając stan teraz,
mówi ona, jak szybko zmienia się każda z nich. To mniej więcej piętnaście
wierszy arytmetyki z sinusami i cosinusami i nie trzeba jej czytać, żeby jej
używać; to \api{double\_pendulum} z biblioteki książki, \pkg{chaoslab}. RK4 z
rozdziału~\ref{ch:motion} stosuje ją krok po kroku, po tysięcznej części
sekundy naraz:

\pyregion{code/ch13/release.py}{setup}{Puszczenie: cztery liczby i długość jednego kroku}{lst:ch13-setup}

\noindent
\code{math.radians} zamienia stopnie na radiany, jednostkę, której oczekują
sinus i cosinus w Pythonie. Pełny obrót to $2\pi$ radianów, nieco ponad sześć,
więc jeden radian to około 57°. Potem listing przesuwa stan naprzód przez
dziesięć sekund i wypisuje go co dwie:

\pyregion{code/ch13/release.py}{run}{Dziesięć sekund wahadła podwójnego, wypisane co dwie}{lst:ch13-run}

\transcript{ch13-release}

\noindent
Najpierw przeczytaj ramiona. Kąt ujemny to wychylenie na drugą stronę, a kąt
powyżej 360° oznacza, że ramię obróciło się dookoła: 758° po dziesięciu
sekundach to dwa pełne obroty i jeszcze 38°. Dolne ramię wiruje. Teraz
przeczytaj ostatnią kolumnę.

\section{Liczba, która nie może drgnąć}
\label{sec:ch13-energy}
\index{energia}\index{symulacja!sprawdzian energii}

\noindent
Wahadło bez tarcia nie może zyskać ani stracić energii: nic go nie popycha i
nic nie spowalnia. Jego energia ma dwie części. Energia wysokości: każdy
ciężarek o masie jednego kilograma wnosi $g \times h$, gdzie $g$ to
\num{9.81}, a wysokość $h$ mierzy się w górę od punktu zawieszenia. Energia
ruchu: połowa kwadratu prędkości każdego ciężarka. Gdy wahadło opada,
wysokość zamienia się w prędkość, i z powrotem; suma pozostaje dokładnie taka
sama.

\begin{think}
Załóżmy, że puszczasz wahadło z obydwoma ramionami ustawionymi poziomo, w tę
samą stronę, pod kątem 90°. Jaka jest jego energia, jeśli wysokości mierzymy
od punktu zawieszenia?
\end{think}
\reveal{Zero. Oba ciężarki są na wysokości punktu zawieszenia, więc obie
wysokości są zerowe, a oba stoją w miejscu, więc nie ma też energii ruchu.}
Wysokości poniżej punktu zawieszenia liczą się jako ujemne, więc wahadło
wiszące w spoczynku ma energię ujemną; znaczenie mają tylko zmiany energii.
Dla puszczenia z wydruku odrobina geometrii umieszcza górny ciężarek pół
metra nad punktem zawieszenia, a dolny cały metr, więc energia wynosi
$\num{9.81} \times (\num{0.5} + 1)$, czyli około \val{ch13.energy.start}
dżula \dash{} od tego zaczyna się kolumna energii. Przeczytaj ją w dół: przez
dziesięć sekund wirowania zmienia się o jedną jednostkę na ósmym miejscu po
przecinku, o stumilionową część dżula.

To jest sprawdzian i to najpożyteczniejszy nawyk w symulowaniu mechaniki.
Komputer nie rozwiązuje reguły; robi małe kroki, każdy odrobinę błędny.
Prawdziwe wahadło zachowuje energię dokładnie, więc każda zmiana energii w
symulacji to zmierzony błąd samych kroków.

\begin{lab}{l13_01_energy}{Liczba, która nie może drgnąć}
Napisz energię wahadła podwójnego sam: \code{heights} podaje wysokość każdego
ciężarka nad punktem zawieszenia, a \code{energy} \dash{} całość. Plik
startowy objaśnia jedyny krok, który wymaga uwagi: dolny ciężarek jest
niesiony przez górny, więc jego prędkość składa się z dwóch ruchów. Test
porównuje twoją energię z energią z biblioteki, a potem obserwuje ją podczas
kołysania, w którym nie może się zmienić.
\end{lab}

\section{Milionowa część radiana}
\label{sec:ch13-apart}
\index{wrażliwość na warunki początkowe!wahadło podwójne}

\noindent
Teraz puść dwa wahadła podwójne. Drugie jest dokładnie takie jak pierwsze, z
wyjątkiem tego, że jego górne ramię zaczyna o jedną milionową część radiana
dalej: na końcu metrowego ramienia to tysięczna część milimetra, dużo mniej
niż grubość włosa. Oba słuchają tej samej dokładnej reguły, z tym samym
krokiem.

\begin{think}
Oba wahadła zachowują energię, jak pokazał sprawdzian, i nic losowego ich nie
dotyka. Czy pod koniec dziesięciu sekund z wydruku powyżej wciąż będą się
kołysać razem? Po minucie? Czy się rozejdą, a jeśli tak, to mniej więcej
kiedy?
\end{think}
\reveal{Rozchodzą się. Po \val{ch13.apart.seconds} sekundy różnią się o więcej
niż jedną dziesiątą radiana, około sześciu stopni, co zobaczyłby każdy
obserwator. Kilka sekund później robią już zupełnie co innego.}

\pyregion{code/ch13/two_releases.py}{gap}{Jak daleko od siebie są dwa wahadła}{lst:ch13-gap}

\pyregion{code/ch13/two_releases.py}{run}{Dwa puszczenia różniące się o milionową część radiana}{lst:ch13-pair}

\transcript{ch13-two-releases}

\noindent
\code{7e-05} to w zapisie Pythona $7 \times 10^{-5}$, czyli siedem
stutysięcznych. Różnica nie rośnie gładko: wznosi się i opada w rytm
kołysania ramion. Ale tendencja jest nieubłagana. Od milionowej do
dziesiątej części to pięć mnożeń przez dziesięć w mniej więcej dziesięć sekund,
więc średnio różnica mnoży się przez dziesięć co dwie sekundy: to efekt
motyla z rozdziału~\ref{ch:butterfly}, w maszynie, którą dałoby się zbudować
na stole warsztatowym.

\plotfig{ch13-two-releases}{Dwa wahadła podwójne puszczone z różnicą
milionowej części radiana. U góry: dolne ramię każdego z nich, nie do
odróżnienia przez około dziesięć sekund. U dołu: różnica między nimi w skali
logarytmicznej, rosnąca nierówno, aż staje się dość duża, by ją
zobaczyć.}{dwa puszczenia, które się rozchodzą}

\noindent
Oto wątpliwość, która powinna ci przyjść do głowy. Każdy krok jest odrobinę
błędny: a co, jeśli rozejście wywołują te błędy, a nie wahadło? Wtedy
dokładniejsze kroki by je przesunęły. Uruchom więc te same dwa puszczenia z
czterema długościami kroku, obserwując przy tym energię:

\pyregion{code/ch13/step_check.py}{parting}{Ta sama para, z dowolną długością kroku}{lst:ch13-parting}

\transcript{ch13-step-check}

\noindent
Każde skrócenie kroku o połowę dzieli błąd energii mniej więcej przez
\val{ch13.energy.shrink}, tak jak obiecywał przepis RK4 z
rozdziału~\ref{ch:motion}, a chwila, w której para się rozchodzi, nie
przesuwa się nawet o setną część sekundy. Przy kroku używanym w książce,
\val{ch13.dt} sekundy, energia zmienia się, zanim para się rozejdzie, o około
\val{ch13.energy.moved} dżula: mniej więcej o jedną część na
\val{ch13.energy.parts}. Rozejście należy do wahadła, nie do komputera.

\plotfig{ch13-energy-error}{Jak daleko energia w symulacji odchodzi od swojej
prawdziwej, stałej wartości, dla czterech długości kroku. Każde skrócenie
kroku o połowę obniża całą krzywą mniej więcej o ten sam czynnik, a linia
kropkowana zaznacza chwilę, w której dwa puszczenia się rozchodzą, taką samą
dla wszystkich czterech.}{sprawdzian energii}

\begin{note}
Ten sprawdzian ma granicę. Policz jedno puszczenie dwa razy, z krokiem
książki i z jego połową, a te dwa obliczenia też się rozejdą, po około
\val{ch13.stepper.apart} sekundach. Drobny błąd każdego kroku jest
pchnięciem, a chaos powiększa je jak każde inne. Zaczyna się jednak od czegoś
dużo mniejszego niż milionowa część radiana, dlatego nasza para rozchodzi się
pierwsza, a każda długość kroku zgadza się co do tego, kiedy.
Rozdział~\ref{ch:computers} pyta, czemu jeszcze można ufać potem.
\end{note}

\begin{trapbox}
\textbf{\enquote{Wahadło, które zachowuje energię, bez tarcia i bez niczego
losowego, nie może być chaotyczne.}} Nie potrzebuje ani tarcia, ani
przypadku. Zachowanie energii tylko częściowo trzyma wahadło w ryzach. Stan
jednego wahadła to dwie liczby; ustalona energia zostawia mu jeden kierunek
ruchu, wokół jednej zamkniętej pętli, i dlatego się powtarza. Stan wahadła
podwójnego to cztery liczby; ustalona energia zostawia trzy kierunki, a trzy
wystarczą do chaosu. Sprawdzian energii dowodzi, że symulacja jest uczciwa.
Nie sprawia, że ruch staje się przewidywalny.
\end{trapbox}

\begin{lab}{l13_02_apart}{Kiedy dwa puszczenia się rozchodzą?}
Zamień listing w funkcję. \code{seconds\_until\_apart} prowadzi krok po kroku
dwa wahadła obok siebie i zwraca pierwszą chwilę, w której różnią się o
więcej niż zadana tolerancja, albo \code{None}, jeśli wciąż są razem, gdy
skończy się czas. Potem jej użyj: o ile dłużej rozchodzą się puszczenia
różniące się o miliardową część radiana niż takie, które różnią się o
milionową? Porównaj odpowiedź z arytmetyką horyzontów z
rozdziału~\ref{ch:lyapunov}.
\end{lab}

\section{Łagodne kołysania i dzikie}
\label{sec:ch13-islands}
\index{twierdzenie KAM}\index{wyspy ruchu regularnego}

\noindent
Wszystko dotąd było puszczeniem z wysoka, z dużą energią. Wahadło podwójne
można też puścić łagodnie.

\begin{think}
Puść to samo wahadło z obydwoma ramionami pod kątem 10° zamiast 120°: ta sama
reguła, ta sama maszyna, to samo pchnięcie o milionową część radiana. Czy
różnica wciąż urośnie w ciągu dwudziestu sekund do czegoś, co każdy by
zobaczył?
\end{think}
\reveal{Nie. Przy 10° oba pozostają w odległości kilku milionowych części
radiana od siebie: pchnięcie prawie nie rośnie. Ta sama maszyna jest spokojna
przy małych wychyleniach i dzika przy dużych.}

\pyregion{code/ch13/gentle_wild.py}{fate}{Co robi para puszczona pod jednym kątem}{lst:ch13-fate}

\transcript{ch13-gentle-wild}

\noindent
Aż do \val{ch13.calm.upto}° każda para jest wciąż razem po
\val{ch13.window} sekundach, a nawet przy \val{ch13.calm.upto}° największa
różnica wynosi \val{ch13.calm.gap} radiana: pchnięcie urosło kilka razy, a
nie sto tysięcy razy. Od \val{ch13.wild.from}° w górę każda para się
rozchodzi. Wykres wypełnia kąty pomiędzy.

\plotfig{ch13-gentle-wild}{Każda kropka to puszczenie obu ramion pod jednym
kątem, razem z kopią przesuniętą o milionową część radiana. Małe wychylenia
pozostają razem przez dwadzieścia sekund; duże się rozchodzą, a zmiana
następuje w przedziale kilku stopni.}{spokojne i dzikie puszczenia}

\noindent
Tak właśnie żyje chaos w układzie bez tarcia. Przy małej energii wahadło
podwójne jest prawie dwoma prostymi kołysaniami, z ramionami poruszającymi
się razem albo przeciwnie, i jego ruch jest regularny. Puść je z większej
wysokości, a pojawia się chaos; z jeszcze większej, a chaos przejmuje władzę,
a ruch regularny przetrwa tylko w zakamarkach: chaotyczne morze z wyspami
ruchu regularnego.

\begin{rigourbox}
Ta książka nie dowodzi, że wahadło podwójne jest chaotyczne; mierzy
wrażliwość dla kilku puszczeń. Twierdzenie stojące za morzem i jego wyspami
nazywa się KAM, od Andrieja Kołmogorowa, który ogłosił je w 1954 roku, oraz
Jürgena Mosera i Władimira Arnolda, którzy udowodnili jego wersje w latach
1962 i 1963. Mówi ono, że gdy regularny układ bez tarcia zostanie lekko
zaburzony, większość jego regularnych ruchów przetrwa, nieco wygięta, a chaos
zostaje zamknięty w cienkich warstwach między nimi; gdy zaburzenie rośnie,
warstwy się poszerzają. Jego dowód to matematyka na poziomie wyższych lat
studiów, do znalezienia w książkach o mechanice takich układów.
\end{rigourbox}

\noindent
Jeszcze jedna różnica względem rozdziału~\ref{ch:lorenz}. Równania Lorenza
zawierają tarcie, a tarcie sprowadza każdy start na jeden atraktor. Wahadło
podwójne nie ma tarcia: każdy start zachowuje swoją energię na zawsze i nie
ma atraktora. Jest za to płot. Energia mówi, które stany w ogóle da się
osiągnąć, a ruch, choćby najdzikszy, pozostaje w środku.

\section{Dwa ciała, trzy ciała}
\label{sec:ch13-three}
\index{problem trzech ciał}\index{Poincaré, Henri}

\noindent
Mechanizmem zegara z rozdziału~\ref{ch:models} był Układ Słoneczny, a jego
triumfem problem dwóch ciał: jednej planety i Słońca. Johannes Kepler odkrył
na początku XVII wieku, że planety poruszają się po elipsach, a Isaac Newton
pokazał w 1687 roku, że jego prawo grawitacji każe dwóm ciałom robić dokładnie
to. Dwa ciała mają wzór: elipsę, zataczaną bez końca. Na każde pytanie o nie
istnieje odpowiedź, którą da się zapisać.

Dodaj trzecie ciało, a wzoru takiego rodzaju w ogóle nie ma. W 1889 roku
Henri Poincaré zdobył nagrodę ufundowaną w imieniu króla Szwecji Oskara II za
pracę nad problemem kilku ciał poruszających się pod wpływem wzajemnej
grawitacji. Gdy jego nagrodzoną rozprawę przygotowywano do druku, znaleziono
w niej błąd. Poprawiając go, Poincaré odkrył, że orbity w pobliżu prostego,
powtarzającego się ruchu trzech ciał mogą owijać się wokół niego w plątaninę
tak zawiłą, że później napisał, iż nawet nie będzie próbował jej narysować.
Poprawiona rozprawa ukazała się w 1890 roku, a to, co opisuje, uznaje się
dziś za pierwszy przebłysk chaosu: regułę równie dokładną jak reguła
Newtona, a za nią ruch, za którym nie nadąży żaden wzór.

W 1908 roku, pisząc dla szerokiej publiczności, Poincaré ujął tę lekcję
słowami, które ta książka właśnie mierzy: przyczyna zbyt mała, by ją
zauważyć, może zdecydować o skutku, którego nie da się przeoczyć, i wtedy
przewidywanie staje się niemożliwe.

\mermaidfig{ch13-bodies}{Dwa ciała mają wzór; trzy nie mają żadnego; planety
są chaotyczne w skali dziesiątek milionów lat, a mimo to pozostają na swoich
miejscach.}{dwa ciała, trzy ciała, planety}

\section{Czy Układ Słoneczny jest zegarem?}
\label{sec:ch13-solar}
\index{Laskar, Jacques}\index{Hyperion}

\noindent
Przez dwa stulecia Układ Słoneczny był wzorem doskonałego zegara. Pod koniec
XVIII wieku Joseph-Louis Lagrange i Pierre-Simon Laplace obliczyli, z
dokładnością, na jaką pozwalały ich metody, że rozmiary orbit planet się
chwieją, ale nie dryfują. Demon Laplace'a z rozdziału~\ref{ch:models} miał
planety za dowód.

W 1989 roku Jacques Laskar obliczył powolne wahania orbit planet na
przestrzeni dwustu milionów lat i stwierdził, że ruch planet wewnętrznych
\dash{} Merkurego, Wenus, Ziemi i Marsa \dash{} jest chaotyczny. Czas
Lapunowa, który znalazł, w sensie z rozdziału~\ref{ch:lyapunov}, wynosił
około pięciu milionów lat: błąd położenia planet mnoży się przez $e$, około
\num{2.7}, co pięć milionów lat.

\begin{think}
Sto milionów lat to dwadzieścia takich czasów Lapunowa. Mniej więcej ile
razy powiększy się błąd w ciągu dwudziestu takich okresów: sto razy, milion
razy czy miliard razy?
\end{think}
\reveal{Około \val{ch13.solar.factor} razy: bliżej miliarda niż miliona.
W tym tempie błąd jednego metra w dzisiejszym położeniu Ziemi urósłby do
\val{ch13.solar.km} kilometrów, więcej niż odległość do Księżyca.}

\begin{trapbox}
\textbf{\enquote{Układ Słoneczny to doskonały zegar.}} W skali kilku tysięcy
lat jest zegarem tak dobrym jak cokolwiek, co mamy: zaćmienia przewiduje się
na stulecia naprzód. W skali dziesiątek milionów lat nie da się go odczytać.
Żaden pomiar i żaden komputer nie powie, gdzie na swojej orbicie będzie
Ziemia za sto milionów lat. Reguła jest dokładna; stan początkowy nie.
\end{trapbox}

\noindent
Nieprzewidywalny to nie to samo co niestabilny. Wynik Laskara nie mówi, że
planety się rozlecą, a późniejsze obliczenia, sięgające miliardów lat, w
niemal każdym przebiegu pokazują planety na orbitach o mniej więcej tych
samych rozmiarach i kształtach. Nie umiemy powiedzieć, gdzie na orbicie
będzie Ziemia za sto milionów lat; umiemy powiedzieć, że najpewniej będzie na
orbicie bardzo podobnej do obecnej. To lekcja z rozdziału~\ref{ch:lorenz} w
nowym miejscu: drogi nie da się przewidzieć, a obszar, po którym się porusza
\dash{} tak.

Chaos ujawnia się szybciej w niektórych zakątkach Układu Słonecznego. Księżyc
Saturna Hyperion to nieforemna bryła na wydłużonej orbicie, a obliczenia
przewidują, że jego obrót nie może się ustalić, tylko koziołkuje
chaotycznie, szarpany przez Saturna raz w jedną, raz w drugą stronę.
Obserwacje z sond Voyager i Cassini zgadzają się z tym przewidywaniem:
Hyperion nie ma stałej doby.

\begin{worldbox}
Geolodzy datują warstwy osadów według rytmów klimatu w nich zapisanych.
Powolne cykle kształtu orbity Ziemi i nachylenia jej osi, nazwane od
nazwiska Milutina Milankovicia, zmieniają to, jak światło słoneczne pada na
planetę, a klimat zostawia w skałach pasujące do nich pasma. Żeby zamienić
pasma w daty, trzeba znać orbitę Ziemi daleko w przeszłość, i geolodzy biorą
ją z obliczeń zespołu Laskara. Tym obliczeniom ufa się, cykl po cyklu, wstecz
do mniej więcej pięćdziesięciu milionów lat i nie dalej \dash{} z powodu
chaosu opisanego w tym podrozdziale.
\end{worldbox}

\begin{summarybox}
\begin{itemize}
\item Stan \textbf{wahadła podwójnego} to cztery liczby, dwa kąty i dwie
  prędkości ich zmian. Jego reguła pochodzi od Newtona, jest dokładna i nie
  ma w niej przypadku, a RK4 stosuje ją krok po kroku.
\item Bez tarcia \textbf{energia} nie może się zmienić, więc energia
  symulacji mierzy błąd kroków: przy kroku książki około jednej części na
  \val{ch13.energy.parts}.
\item Puszczenia różniące się o milionową część radiana rozchodzą się po
  \val{ch13.apart.seconds} sekundy przy każdej długości kroku:
  \textbf{rozejście to fizyka}. Zachowanie energii nie zapobiega chaosowi.
\item Łagodne puszczenia są \textbf{regularne}, a dzikie chaotyczne: morze z
  wyspami, jak opisuje twierdzenie KAM. Bez tarcia nie ma atraktora; energia
  ogradza ruch.
\item Dwa ciała mają wzór, a trzy nie. Poprawiona nagrodzona rozprawa
  \textbf{Poincarégo} z 1890 roku była pierwszym przebłyskiem chaosu.
\item Laskar stwierdził, że planety wewnętrzne są chaotyczne, a ich
  \textbf{czas Lapunowa} wynosi około pięciu milionów lat: ich położeń nie da
  się przewidzieć na sto milionów lat naprzód, choć ich orbity pozostają
  mniej więcej takie, jakie są.
\end{itemize}
\end{summarybox}

\canyou

\begin{checkpoint}
\item Dlaczego wahadło podwójne może być chaotyczne, a pojedyncze nie, choć
  żadne z nich nie traci energii?
  \answerto{Ustalenie energii dwuliczbowego stanu pojedynczego wahadła
  zostawia jeden kierunek, wokół zamkniętej pętli. Ustalenie jej dla czterech
  liczb wahadła podwójnego zostawia trzy, a to wystarcza do chaosu.}
\item Wahadło podwójne puszczono ze spoczynku z obydwoma ramionami
  skierowanymi pionowo w górę. Jaka jest jego energia, jeśli wysokości
  mierzymy od punktu zawieszenia?
  \answerto{Ciężarki są metr i dwa metry wyżej, więc $\num{9.81} \times
  (1 + 2)$, czyli około \num{29.4} dżula.}
\item Twoja symulacja maszyny bez tarcia traci w ciągu minuty jedną dziesiątą
  energii. Co z tego wnioskujesz?
  \answerto{Błąd kroków jest za duży, by ufać wynikowi. Skracaj krok, aż
  energia się utrzyma, a potem sprawdź, czy twoja odpowiedź przetrwa kolejne
  skrócenie kroku o połowę.}
\item Dwa puszczenia rozchodzą się po dziesięciu sekundach. Skracasz krok o
  połowę i teraz rozchodzą się po trzech. Co ci to mówi?
  \answerto{Że rozejście wywołały przynajmniej częściowo błędy kroków, więc
  żadnej z tych liczb nie można ufać. Rozejście, które jest fizyką, nie
  przesuwa się, gdy zmienia się krok.}
\item Wahadło podwójne nie ma atraktora. Co powstrzymuje jego ruch przed
  pójściem gdziekolwiek?
  \answerto{Jego energia, która nigdy się nie zmienia i decyduje, które stany
  w ogóle da się osiągnąć.}
\item Przy czasie Lapunowa wynoszącym pięć milionów lat mniej więcej przez
  ile mnoży się błąd w ciągu dziesięciu milionów?
  \answerto{Dwa razy przez $e$: $\num{2.7} \times \num{2.7}$, nieco ponad
  siedem.}
\item Czy chaos Układu Słonecznego oznacza, że planety się rozlecą?
  \answerto{Nie. Ich położeń nie da się przewidzieć dalej niż na kilkadziesiąt
  milionów lat, ale obliczenia sięgające miliardów lat w niemal każdym
  przebiegu pokazują, że ich orbity zachowują mniej więcej te same rozmiary i
  kształty.}
\end{checkpoint}
